\documentclass[11pt]{amsart}


\usepackage{amscd, amssymb}
\usepackage{enumerate}
\usepackage[all]{xypic}
\usepackage{multirow}
\usepackage{hhline}
\usepackage{verbatim}
\usepackage{mathdots}
\usepackage{comment}
\usepackage{tikz-cd}
\usepackage[titletoc]{appendix}
\usepackage{stmaryrd}%this conflicts with extpfeil
\usepackage{graphicx}
\usepackage{mathtools}
%\usepackage{extpfeil}
%\usepackage{amsrefs}
%\usepackage{thmtools}
%\usepackage{float}
%\usepackage{makeidx}

%\usepackage{amsthm}
\usepackage{leftidx}
\usepackage{mathabx }
\usepackage{amsfonts}
\usepackage{mathrsfs}

\usepackage{dynkin-diagrams}
\usepackage{xcolor}

\usepackage{enumitem}
%\usepackage{enumerate}

\usepackage[margin=0.8in]{geometry}


\usepackage[colorlinks=true, linkcolor=blue,
 citecolor=blue, urlcolor=blue]{hyperref}

%\usetikzlibrary{decorations.markings}


%\renewcommand{\thesection}{\Roman{section}}

%%% table of contents %%%
% \setcounter{tocdepth}{2}
% \let\oldtocsection=\tocsection
% \renewcommand{\tocsection}[2]{\hspace{0em}\oldtocsection{#1}{#2}}





\theoremstyle{plain}
\newtheorem{Thm}[equation]{Theorem}
\newtheorem*{Thm*}{Theorem}
\newtheorem{Cor}[equation]{Corollary}
\newtheorem{Prop}[equation]{Proposition}
\newtheorem{Lem}[equation]{Lemma}


\theoremstyle{definition}
\newtheorem*{Rmk}{Remark}
\newtheorem*{Rmks}{Remarks}
\newtheorem{Def}[equation]{Definition}
\newtheorem*{hypo}{Working
Hypotheses}


\numberwithin{equation}{section}



%\newcommand{\A}{\mathbb{A}}
\newcommand{\C}{\mathbb{C}}
\newcommand{\F}{\mathbb{F}}
%\newcommand{\HH}{\mathbb{H}}
%\newcommand{\K}{\mathbb{K}}
\newcommand{\N}{\mathbb{N}}
\newcommand{\Q}{\mathbb{Q}}
\newcommand{\R}{\mathbb{R}}
%\newcommand{\V}{\mathbb{V}}
%\newcommand{\W}{\mathbb{W}}
%\newcommand{\X}{\mathbb{X}}
%\newcommand{\Y}{\mathbb{Y}}
\newcommand{\Z}{\mathbb{Z}}
%\newcommand{\OOmega}{\mathbf{\Omega}}



\newcommand{\OF}{\mathcal{O}_F}
\newcommand{\OFd}{\mathcal{O}_F^{\diamondsuit}}


%\newcommand{\Gm}{\mathbb{G}_m}
\newcommand{\GL}{\operatorname{GL}}
%\newcommand{\PGL}{\operatorname{PGL}}
\newcommand{\SL}{\operatorname{SL}}
\newcommand{\Sp}{\operatorname{Sp}}
%\newcommand{\GSp}{\operatorname{GSp}}
\newcommand{\SO}{\operatorname{SO}}
\newcommand{\OO}{\operatorname{O}}
%\newcommand{\UU}{\operatorname{U}}
%\newcommand{\SU}{\operatorname{SU}}
\newcommand{\Mp}{\operatorname{Mp}}

\newcommand{\Rad}{\operatorname{rad}}

\newcommand{\pr}{{\operatorname{pr}}}



\newcommand{\xa}{x_{\alpha}}
\newcommand{\wa}{w_{\alpha}}
\newcommand{\ha}{h_{\alpha}}
\newcommand{\xb}{x_{\beta}}
\newcommand{\wb}{w_{\beta}}
\newcommand{\hb}{h_{\beta}}

\newcommand{\xta}{\tilde{x}_{\alpha}}
\newcommand{\wta}{\tilde{w}_{\alpha}}
\newcommand{\hta}{\tilde{h}_{\alpha}}
\newcommand{\xtb}{\tilde{x}_{\beta}}
\newcommand{\wtb}{\tilde{w}_{\beta}}
\newcommand{\htb}{\tilde{h}_{\beta}}



\newcommand{\alphac}{{\check{\alpha}}}
\newcommand{\betac}{{\check{\beta}}}
\newcommand{\gammac}{{\check{\gamma}}}

\newcommand{\g}{\mathfrak{G}}
\newcommand{\h}{\mathfrak{h}}
\renewcommand{\sp}{\mathfrak{sp}}


\newcommand{\Spt}{\widetilde{\operatorname{Sp}}}
\newcommand{\SLt}{\widetilde{\operatorname{SL}}}
\newcommand{\Bt}{\overline{B}}
\newcommand{\Mt}{\widetilde{M}}
\newcommand{\Nt}{\widetilde{N}}
\newcommand{\Pt}{\widetilde{P}}
\newcommand{\Qt}{\widetilde{Q}}
\newcommand{\Gt}{\overline{G}}
\newcommand{\Ht}{\overline{H}}
\newcommand{\It}{\overline{I}}
\newcommand{\Jt}{\widetilde{J}}
\newcommand{\Kt}{\overline{K}}
\newcommand{\Tt}{\overline{T}}
\newcommand{\Xt}{X_{Q,n}}
\newcommand{\XXt}{\widetilde{\mathbb{X}}}
\newcommand{\Yt}{Y_{Q,n}}
\newcommand{\Wt}{\widetilde{W}}
\newcommand{\Omegat}{\widetilde{\Omega}}

\newcommand{\Gbbt}{\overline{{\Gbb}}}
\newcommand{\Tbbt}{\overline{{\Tbb}}}
\newcommand{\Bbbt}{\overline{{\Bbbb}}}

\newcommand{\Gbb}{{\mathbb{G}}}
\newcommand{\Tbb}{{\mathbb{T}}}
\newcommand{\Bbbb}{{\mathbb{B}}}

\newcommand{\Gammat}{\widetilde{\Gamma}}
\newcommand{\Gammaft}{\widetilde{\mathsf{\Gamma}}}
\newcommand{\Gammatp}{\widetilde{\Gamma}'}


\newcommand{\mt}{\tilde{m}}
\newcommand{\gt}{\tilde{g}}
\newcommand{\hti}{\tilde{h}}
\newcommand{\kt}{\tilde{k}}
\newcommand{\pt}{\tilde{p}}
\newcommand{\qt}{\tilde{q}}
\newcommand{\ttt}{\overline{t}}

\newcommand{\xt}{\tilde{x}}
\newcommand{\yt}{\overline{y}}
\newcommand{\zt}{\tilde{z}}
\newcommand{\wt}{\tilde{w}}
\newcommand{\htt}{\tilde{h}}

\newcommand{\yh}{\hat{y}}

\newcommand{\taut}{\tilde{\tau}}
\newcommand{\pit}{\tilde{\pi}}


\newcommand{\nuh}{\hat{\nu}}


\newcommand{\tbb}{\widebar{t}}
\newcommand{\wbb}{\widebar{w}}
\newcommand{\xbb}{\widebar{x}}
\newcommand{\ybb}{\widebar{y}}
\newcommand{\zbb}{\widebar{z}}





\newcommand{\Ind}{\operatorname{Ind}}
\newcommand{\cInd}{\operatorname{c-Ind}}
\newcommand{\ind}{\operatorname{ind}}
\newcommand{\Irr}{\operatorname{Irr}}
\newcommand{\diag}{\operatorname{diag}}
\newcommand{\supp}{\operatorname{supp}}
\newcommand{\ch}{\operatorname{ch}}
\newcommand{\Span}{\operatorname{span}}
\newcommand{\rank}{\operatorname{rank}}
\newcommand{\End}{\operatorname{End}}
\newcommand{\Hom}{\operatorname{Hom}}
\newcommand{\Aut}{\operatorname{Aut}}
\newcommand{\ord}{\operatorname{ord}}
\newcommand{\vol}{\operatorname{vol}}
\newcommand{\sign}{\operatorname{sign}}
\newcommand{\sig}{\operatorname{sig}}
\newcommand{\Herm}{\operatorname{Herm}}
\newcommand{\Quad}{\operatorname{Quad}}
\newcommand{\Sym}{\operatorname{Sym}}
\newcommand{\image}{\operatorname{Im}}
\newcommand{\Int}{\operatorname{Int}}
\newcommand{\Gal}{\operatorname{Gal}}
\newcommand{\Res}{\operatorname{Res}}
\newcommand{\disc}{\operatorname{disc}}
\newcommand{\tr}{\operatorname{tr}}
\newcommand{\Tr}{\operatorname{Tr}}
\newcommand{\rad}{\operatorname{rad}}

\newcommand{\Aff}{\operatorname{Aff}}
\newcommand{\aff}{\mathrm{aff}}
\newcommand{\af}{\mathrm{af}}
\newcommand{\ea}{\mathrm{ea}}



\newcommand{\Sm}{{\mathcal{S}}}

\newcommand{\ee}{\mathbf{e}}
\newcommand{\ff}{\mathbf{f}}
\newcommand{\rr}{\mathbf{r}}
\newcommand{\hh}{\mathbf{h}}
\newcommand{\xx}{\mathbf{x}}
\newcommand{\yy}{\mathbf{y}}
\newcommand{\ww}{\mathbf{w}}


\newcommand{\s}{\mathbf{s}}

\newcommand{\rp}{\mathbf{r}_{\psi}}

\newcommand{\Aa}{A^{(2)}}
\newcommand{\Bb}{B^{(2)}}
\newcommand{\II}{\mathcal{I}}
\newcommand{\JJ}{\mathcal{J}}
\newcommand{\KK}{\mathcal{K}}
\newcommand{\LL}{\mathcal{L}}
\newcommand{\MM}{\mathcal{M}}

\newcommand{\KKt}{\widetilde{\mathcal{K}}}
\newcommand{\IIt}{\widetilde{\mathcal{I}}}
\newcommand{\JJt}{\widetilde{\mathcal{J}}}

\newcommand{\X}{\mathfrak{X}}


\newcommand{\Scal}{\mathcal{S}}
\newcommand{\Ocal}{\mathcal{O}}
\newcommand{\Fcal}{\mathcal{F}}
\newcommand{\Wcal}{\mathcal{W}}
\newcommand{\Hcal}{\mathcal{H}}
\newcommand{\Ncal}{\mathcal{N}}
\newcommand{\Dcal}{\mathcal{D}}


\newcommand{\Bf}{\mathsf{B}}
\newcommand{\Ff}{\mathsf{F}}
\newcommand{\Hf}{\mathsf{H}}
\newcommand{\If}{\mathsf{I}}
\newcommand{\Kf}{\mathsf{K}}
\newcommand{\Mf}{\mathsf{M}}
\newcommand{\Nf}{\mathsf{N}}
\newcommand{\Tf}{\mathsf{T}}
\newcommand{\Uf}{\mathsf{U}}
\newcommand{\Vf}{\mathsf{V}}
\newcommand{\Gammaf}{\mathsf{\Gamma}}

\newcommand{\Hft}{\widetilde{\mathsf{H}}}
\newcommand{\Ift}{\widetilde{\mathsf{I}}}
\newcommand{\Kft}{\widetilde{\mathsf{K}}}
\newcommand{\Mft}{\widetilde{\mathsf{M}}}
\newcommand{\Nft}{\widetilde{\mathsf{N}}}
\newcommand{\Bft}{\widetilde{\mathsf{B}}}
\newcommand{\Pft}{\widetilde{\mathsf{P}}}
\newcommand{\Tft}{\widetilde{\mathsf{T}}}



\newcommand{\la}{\langle}
\newcommand{\ra}{\rangle}
\newcommand{\lla}{\langle\!\langle}
\newcommand{\rra}{\rangle\!\rangle}
\newcommand{\llp}{(\!(}
\newcommand{\rrp}{)\!)}

\newcommand{\ie}{{\it i.e.}\ }
%\newcommand{\apriori}{{\it a priori }}
%\newcommand{\afortiori}{{\it a fortiori }}

\newcommand{\st}{\,:\,}
\newcommand{\qand}{{\quad\text{and}\quad}}
\newcommand{\qqand}{{\qquad\text{and}\qquad}}
\newcommand{\qor}{{\quad\text{or}\quad}}
\newcommand{\by}{{\quad\text{by}\quad}}

\newcommand{\ve}{\varepsilon}


\newcommand{\one}{\mathbf{1}}


\title[High degree covers and genericity]{High degree covers and genericity}


\author{Edmund Karasiewicz}

\subjclass[2020]{11F70, 22E50}
\keywords{p-adic groups; Brylinski–Deligne covers; types; Whittaker funcitonals}





% \setlength{\oddsidemargin}{0.2in} \setlength{\evensidemargin}{0.2in}
%\setlength{\textwidth}{5.8in}


%\pgfkeys{/Dynkin diagram, edge-length=1cm, root-radius=.1cm}



\def\shui#1{\textcolor{blue}{#1}}

\allowdisplaybreaks


\begin{document}


\begin{abstract}
    It is known that the genuine irreducible Iwahori-spherical representations of a tame $n$-fold Brylinski-Deligne cover of a $p$-adic split connected reductive group are all generic, when the degree $n$ is suitably large. This has been established by studying the the structure of the Gelfand-Graev representation as a module under the action of the Iwahori-Hecke algebra.  
    
    In this paper, we present a different proof based on metaplectic Moy-Prasad theory. 
\end{abstract}

\maketitle


\setcounter{tocdepth}{1}
%\tableofcontents

%%%%%%%%%%%%%%%%%%%%%%%%%%%%%%%%%%%%%%%%%%%%%%%%%%%%%%%%%%%%%%%%%%%%%%%

\section{Introduction}

%%%%%%%%%%%%%%%%%%%%%%%%%%%%%%%%%%%%%%%%%%%%%%%%%%%%%%%%%%%%%%%%%%%%%%%  

In the representation theory of tame $n$-fold Brylinski--Deligne (BD) covers of split connected reductive $p$-adic groups, it is known that genuine irreducible representation in the Iwahori-spherical block (equivalently the block of unramified principal series \cite{S04}) are all generic, if $n$ is suitably large. This was proved in \cite{GGK24,GGK25} by examining the Gelfand-Graev representation as a module over a Hecke algebra. Earlier work on the genericity of theta representations (also called exceptional representations) appears in Kazhdan--Patterson \cite{KP84} for Kazhdan-Patterson covers and Gao \cite{Gao_theta_17} for more general BD covers.

The objective of this note is to present a simple explanation for this phenomenon based on Moy--Prasad theory \cite{MP94,MP96}, which was extended to tame covers by Howard--Weissman \cite{HW09}. 

We note that Moy--Prasad theory already plays a role in recent calculations of the stable wave front set of certain representations in the Iwahori-block of tame BD covers \cite{KOW26,Gao-Wang26}. However, these works involve the local character expansion and as a result impose constrains on the residual characteristic beyond the tame assumption, $p$ does not divide $n$. Here, we only assume that $p$ does not divide $n$. 


%The more general question of computing the (stable) wave front set of representations in the Iwahori-block was recently considered in \cite{Gao-Wang26}.

%(See \cite{Gao_theta_17} for an extension to BD covers.) The result in the more general setting above is proved in \cite{GGK24} and further examined in the follow-up \cite{GGK25}. 



%See also \cite{Gao-Wang26}, for calculations of the wave front set of certain Iwahori-spherical representations.

Let $F$ be a finite extension of $\Q_{p}$. Let $\Ocal\subset F$ be the ring of integers and fix $\varpi\in\Ocal$ a uniformizer. Let $\kappa$ be the residue field of $F$ and $q$ its cardinality. 

For a field $L$ and $m\in\Z_{\geq 1}$, we will write $\mu_{m}(L)$ for the group of $m$-th roots of unity in $L$. For the $p$-adic field $F$ we set $\mu_{m}:=\mu_{m}(F)$. 

%Let $\Gbb$ be a split connected reductive linear algebraic group over $F$. 

Let $G$ be the group of $F$-points of a split connected reductive linear algebraic group over $F$. Let $T\subset G$ be a maximal split torus with character lattice $X$ and cocharacter lattice $Y$. Let $\Phi\subset X$ be the subset of roots and $\Phi^{\vee}\subset Y$ the subset of corroots attached to the pair $(G,T)$. For $\alpha\in\Phi$ we write $U_{\alpha}$ for the root subgroup attached to $\alpha$.

Let $B=B^{+}\supset T$ be a Borel subgroup with unipotent radical $U$. Let $\Delta\subset\Phi$ be the set of simple roots determined by $B$. Let $B^{-}\supset U^{-}$ be the opposite Borel and unipotent radical.

Let $K$ be a hyperspecial maximal compact subgroup such that $T_{0}:=T\cap K$ is the maximal compact subgroup of $T$. Let $I^{\pm}\subset K$ be an Iwahori subgroup such that $B^{\pm}\cap K = B^{\pm}\cap I$. For a pro-finite group $H$, let $H_{+}$ be its maximum pro-$p$ subgroup, when it exists. We will write $I^{-}_{+}:=(I^{-})_{+}$.

Assume that $|\mu_{n}|=n$. Let $\Gt$ be a topological central extension of $G$ by $\mu_{n}$ arising from the construction of Brylinski-Deligne \cite{BD01}. The group $\Gt$ fits into a central extension of the form
\begin{equation}\label{Eqn:CExt}
    \mu_{n}\hookrightarrow \Gt\stackrel{\mathrm{pr}}{\twoheadrightarrow} G.
\end{equation}
For any $H\subset G$, let $\overline{H}:=\mathrm{pr}^{-1}(H)$. Let $\epsilon=\epsilon_{n}:\mu_{n}\rightarrow \C^{\times}$ be a faithful character. A $\Gt$-representation is called $\epsilon$-genuine if $\mu_{n}$ acts through $\epsilon$. Let $\mathrm{Irr}_{\epsilon}(\Gt)$ be the set of isomorphism classes of $\epsilon$-genuine irreducible smooth $\Gt$-representations.

Recall that the covering group $\Gt$ comes with $Q:Y\rightarrow \Z$ a $W$-invariant quadratic form, where $W$ is the Weyl group of $T$. Let $B=B_{Q}$ be the bilinear form on $Y$ associated to the quadratic form $Q$. For any $m\in\Z$, let $Y_{Q,m}:=\{y\in Y|B(y,z)\in m\Z\text{, for all }z\in Y\}$.

For the remainder of this note we fix $n\in\Z_{\geq 1}$ such that $\gcd(p,n)=1$ and $|\mu_{n}|=n$. This implies that $n$ divides $q-1$. In particular, $\Gt$ is a tame BD cover.

We will also assume that there exists $\s_{K}:K\rightarrow \Kt$, a splitting of the sequence \eqref{Eqn:CExt} over $K$. This happens, for example, when $\Gt$ is incarnated by $(D,\eta)$, where $D$ is a fair bisector and the group homomorphism $\eta:\Z[\Phi^{\vee}]\rightarrow F^{\times}$ is valued in $\Ocal^{\times}F^{\times n}$ \cite[Theorem 4.2 and subsequent remarks]{GG18}.

Now we state our main theorem.

\begin{Thm}\label{Thm:Main_Thm}
    Let $\Gt$ be an $n$-fold BD cover of a split connected reductive group $G$ over a $p$-adic field $F$ such that $|\mu_{n}|=n$, $\gcd(p,n)=1$, and $K$ splits \eqref{Eqn:CExt} with splitting $\s_{K}$. Let $\pi\in \mathrm{Irr}_{\epsilon}(\Gt)$. Let $N$ be the number of free $W$-orbits in $Y/\Yt$. 

    If $\pi^{\s_{K}(I)}\neq 0$, then the Whittaker dimension with respect to any non-degenerate $\psi\in \mathrm{Hom}(U^{-},\C^{\times})$ of conductor $(\varpi)$ is at least $N$. 
\end{Thm}


\begin{Rmk}
    Earlier we used the informal statement `$n$ is suitably large' (this also hides an implicit dependence on $Q$). In the theorem this informal statement is replaced by `$Y/\Yt$ contains a free $W$-orbit'. 
\end{Rmk}

Actually, we will prove the more general Theorem \ref{Thm:generalized_main}, from which Theorem \ref{Thm:Main_Thm} follows as a corollary.

Our main theorem is the result of two inputs. The first input follows from unraveling some results from depth $0$ metaplectic type theory. The main consequence of this is that in the setting of Theorem \ref{Thm:Main_Thm}, the $\Gt$-representations contain regular principal series among their $\Kf$-types, which are generic. The second input is a result lifting Whittaker functionals from finite groups to $p$-adic groups and their covers.

\section{Metaplectic Moy-Prasad theory}

Moy-Prasad theory was developed in \cite{MP94,MP94}, in the setting of linear groups. These ideas were generalized to tame covering groups by Howard--Weissman \cite{HW09}. We will specialize some results of Moy--Prasad and Howard--Weissman and then examine their effect on the existence of Whittaker functionals. 

The main point for this paper is that the set of genuine depth $0$ Iwahori types admits a free action of the finite abelian group $Y/\Yt$, and this proliferation of distinct types can create many independent Whittaker functionals when the degree of the cover $n$ is suitably large. 

We begin with the commutator of $\Tt$, which is a source of non-genuine characters. For $y\in Y$, let $\phi_{y}\in \mathrm{Hom}(T,\mu_{n}(\C))$ be the character defined as follows. Let $t\in T$ and $y\in Y$. Fix a uniformizer $\varpi\in F^{\times}$. Let $\ttt\in \mathrm{pr}^{-1}(t)$ and $\yt(\varpi)\in \mathrm{pr}^{-1}(y(\varpi))$ be arbitrary. Let $[-,-]$ be the commutator in $\Tt$. We define
\begin{equation*}
    \phi_{y}(t):= \epsilon([\ttt,\yt(\varpi)]).
\end{equation*}
Under the identification $T\cong Y\otimes F^{\times}$, for $z\in Y$, $a\in F^{\times}$ we have
\begin{equation}\label{Eq:Tt_Comm_Form}
    \phi_{y}(z\otimes a) = (\varpi,a)_{n}^{B(y,z)},
\end{equation}
where $(-,-)_{n}$ is the $n$-th order Hilbert symbol for $F$.

The next lemma follows from well-known properties of the tame symbol $(-,-)_{n}$ and identity \eqref{Eq:Tt_Comm_Form}. 

\begin{Lem}\label{Lem:phi_y_props}
    Assume that $p$ does not divide $n$. Then:
    \begin{enumerate}
        \item\label{Lem:phi_y_props:1} $\phi_{y}|_{T_{0}}$ does not depend on the choice of $\varpi$;
        \item\label{Lem:phi_y_props:2} $\phi_{y}$ is trivial on $T_{1}:=T\cap I_{+}$, the maximal pro-$p$ subgroup of $T$.
        \item\label{Lem:phi_y_props:3} $\phi_{y}|_{T_{0}}$ is trivial if and only if $y\in \Yt$.
    \end{enumerate}
\end{Lem}

By Lemma \ref{Lem:phi_y_props} Item \eqref{Lem:phi_y_props:2}, $\phi_{y}|_{T_{0}}$ factors through $T_{\kappa}:=T_{0}/T_{1}$. 

\begin{Lem}\label{Lem:Char_equi}
    The map $Y/\Yt\hookrightarrow\mathrm{Hom}(T_{\kappa},\mu_{n}(\C))$ defined by $y\mapsto \phi_{y}|_{T_{0}}$ is a $W$-equivariant injective group homomorphism.

    %Moreover, if $\Yt=nY$, then the map is a bijection.
\end{Lem}

\begin{proof}
    One can directly check that the map is a $W$-equivariant group homomorphism. By Lemma \ref{Lem:phi_y_props} Item \eqref{Lem:phi_y_props:3}, the character $\phi_{y}$ is trivial if and only if $y\in \Yt$. Thus the map is injective.
\end{proof}

The inclusion $T_{0}\subset I$ induces an isomorphism $T_{\kappa}\cong I/I_{+}$. Thus we can also view $\phi_{y}|_{T_{0}}$ as a character of $I$. From now on we will just write $\phi_{y}$, instead of $\phi_{y}|_{T_{0}}$, for this character on $T_{0}$, $T_{\kappa}$, or $I$.

Let $\s_{I}:=\s_{K}|_{I}$. Then $\s_{I}$ is a splitting of $I$. Similarly, $\s_{I_{+}}:=\s_{K}|_{I_{+}}$ is a splitting of $I_{+}$. 

Since we are assuming that $\gcd(p,n)=1$, it follows that $I_{+}$ has a unique splitting. In contrast, the splitting of $I$ is not unique; every splitting of $I$ is of the form $\phi\cdot \s_{I}$, where $\phi\in\mathrm{Hom}(I,\mu_{n})\cong \mathrm{Hom}(T_{\kappa},\mu_{n})$. 

%Note that the character $\phi_{\s_{I}}|_{\s_{I}(I)}$ is by definition the trivial character. The pairs $(\It, \phi_{y}\cdot \phi_{\s_{I}})$, where $y\in Y$, are $\epsilon$-genuine depth $0$ minimal $\mathsf{K}$-types.

The next simple lemma highlights a special feature of the metaplectic Iwahori type theory. Recall that the important notion of associate types was introduced by Moy-Prasad \cite{MP94}. For the setting of tame covers see \cite[Definition 3.4]{HW09}. 

\begin{Prop}\label{Prop:Y_type_action}
    Let $y\in Y$ and $w\in W$. Let $\chi$ be any $\epsilon$-genuine depth $0$ character of $\It$.
    
    The two depth $0$ types $(\It,\chi)$ and $(\It,\phi_{y}\cdot \,^{w}\chi)$ are associate in $\Gt$.
\end{Prop}

\begin{proof}
    Let $\yt(\varpi)\in \mathrm{pr}^{-1}(y(\varpi))$ be arbitrary and let $k\in \Kt\cap N_{\Gt}(\Tt)$ be any representative of $w^{-1}\in W$. The result follows because for any $t\in \Tt_{0}$ the identity \eqref{Eq:Tt_Comm_Form} implies that $\chi(k\yt(\varpi)t\yt(\varpi)^{-1}k^{-1}) = (\phi_{y}\cdot \,^{w}\chi)(t)$.
\end{proof}

\begin{Rmk}
    From Lemma \ref{Lem:Char_equi} and Proposition \ref{Prop:Y_type_action} we get an action of the extended affine Weyl group $W\ltimes Y$, factoring through the finite group $W\ltimes Y/\Yt$, on the $\epsilon$-genuine depth $0$ Iwahori types. Moreover, the action of $Y/\Yt$ is free.
\end{Rmk}

The next corollary follows by combining Proposition \ref{Prop:Y_type_action} and specializing \cite[Proposition 3.7]{HW09}, which is the tame covering analog of \cite[Proposition 6.2]{MP96}.

\begin{Cor}\label{Cor:Translate_Types}
    Let $\pi\in\mathrm{Irr}_{gen}(\Gt)$. Let $\chi:\It\rightarrow \C^{\times}$ be a depth $0$ $\epsilon$-genuine character of $\It$.
    
    If $\pi^{(\It,\chi)}\neq 0$, then $\pi^{(\It,\phi_{y}\cdot \,^{w}\chi)}\neq 0$ for all $y\in Y$ and $w\in W$.
\end{Cor}

In the next proposition we describe how the appearance of the depth $0$ Iwahori type $(\s_{I}(I),\chi)$ and some root system combinatorics translates into the setting of finite groups of Lie type. Let $G_{\kappa}= K/K_{+}$, $B_{\kappa} = I/K_{+}$, and $U^{-}_{\kappa} = I^{-}_{+}/K_{+}$. Since $T_{\kappa}\cong I/I_{+}$, the natural quotient map gives a surjective group homomorphism $B_{\kappa}\rightarrow T_{\kappa}$.

\begin{Prop}\label{Prop:Finite_Generic}
    Let $\pi\in\mathrm{Irr}_{\epsilon}(\Gt)$. Let $\psi_{\kappa}$ be a non-degenerate character of $U^{-}_{\kappa}$. Let $\chi\in\mathrm{Hom}(T_{\kappa},\C^{\times})$, which can be inflated to a character of $\s_{I}(I)$. Let $\Ocal_{\chi}\subset \mathrm{Hom}(T_{\kappa},\C^{\times})$ be the $W\ltimes Y/\Yt$-orbit of $\chi$. 

    Let $N$ be equal to the number of free $W$-orbits in $\Ocal_{\chi}$.

    If $\pi^{(\s_{I}(I),\chi)}\neq 0$, then the $G_{\kappa}$-representation $\pi^{K_{+}}$ contains at least $N$ non-isomorphic $\psi_{\kappa}$-generic irreducible constituents. More specifically, these $N$ constituents can be taken to be regular principal series.
\end{Prop}

\begin{proof}
    By Proposition \ref{Cor:Translate_Types}, $\pi^{(\s_{I}(I),\chi^{\prime})}\neq 0$ for all $\chi^{\prime}\in \Ocal_{\chi}$. Let $\chi_{reg}\in \Ocal_{\chi}$ be any character such that the $W$-orbit of $\chi_{reg}$ is free, if it exists. Since the regular principal series $\mathrm{Ind}_{B_{\kappa}}^{G_{\kappa}}\chi_{reg}$ is irreducible, Frobenius reciprocity implies that  $\mathrm{Ind}_{B_{\kappa}}^{G_{\kappa}}\chi_{reg}$ is a $G_{\kappa}$-submodule of $\pi^{K_{+}}$. We also know that $\mathrm{Ind}_{B_{\kappa}}^{G_{\kappa}}\chi_{reg}$ is generic for any non-degenerate character. %Thus $\pi^{K_{+}}$ contains a $\psi_{\kappa}$-generic constituent. 

    Finally, since distinct free $W$-orbits in $\Ocal_{\chi}$ produce nonisomorphic regular principal series in $\pi^{K_{+}}$, it follows that the number of distinct $\psi_{\kappa}$-generic constituents is at least $N$.
\end{proof}

\section{Lifting Whittaker functionals}

In this section we prove Proposition \ref{Prop:Whit_lift}, which essentially relates Whittaker functionals on $\pi^{K_{+}}$ to Whittaker functionals on $\pi$. We suspect that this result is known, but we do not know a reference.

The sequence \eqref{Eqn:CExt} splits uniquely over any (algebraic) unipotent subgroup. We use this splitting to identify $U^{\pm}$ with its image in $\Gt$. Let $\psi\in\mathrm{Hom}(U^{-},\C^{\times})$ be a nondegenerate character of conductor $(\varpi)$. This means for every $\alpha\in\Delta$ the character $\psi$ is trivial on $U_{-\alpha,1}:=U_{-\alpha}\cap K_{+}$ and nontrivial on $U_{-\alpha,0}:=U_{-\alpha}\cap K$.

The inclusion $U^{-}\cap K\subset I^{-}_{+}$ descends to an isomorphism $(U^{-}\cap K)/(U^{-}\cap K_{+})\cong U^{-}_{\kappa}$. Since $\psi$ has conductor $(\varpi)$, the character $\psi|_{U^{-}\cap K}$ descends to a non-degenerate character of $(U^{-}\cap K)/(U^{-}\cap K_{+})$. Thus we may view $\psi|_{U^{-}\cap K}$ as a character of $U^{-}_{\kappa}$, which we call $\psi_{\kappa}$. We also write $\psi_{\kappa}$ for this character pulled back to $I^{-}_{+}$.

We begin with some notation. Let $\mathrm{Wh}_{\psi}(\pi):=\mathrm{Hom}_{G}(\pi,\mathrm{Ind}_{\mu_{n}U^{-}}^{G}\epsilon\otimes\psi)$ be the space of Whittaker models and $\mathcal{V}:=\mathrm{ind}_{\mu_{n}U^{-}}^{G}(\epsilon^{-1}\otimes\psi^{-1})$ the Gelfand-Graev representation. Let $\mathcal{W}:=\mathrm{ind}_{\mu_{n}I^{-}_{+}}^{G}\epsilon^{-1}\otimes\psi_{\kappa}^{-1}$. Let $\Hcal$ be the convolution algebra of locally constant compactly supported functions $f:\Gt\rightarrow \C$ such that $f(\zeta\gamma_{1}g\gamma_{2})=\epsilon(\zeta)^{-1}f(g)$, where $\zeta\in\mu_{n}$ and $\gamma_{j}\in \s_{I_{+}}(I_{+})$.

Recall that $\mathcal{W}^{\s_{I_{+}}(I_{+})}\cong \mathcal{V}^{\s_{I_{+}}(I_{+})}$ as $\Hcal$-modules, by \cite[Lemmas 5.1-5.3]{GGK24}. (This isomorphism was brought to our attention by the referee of \cite{GGK24}.) 

For $\pi$ a smooth admissible $\Gt$-representation, we write $\pi^{\vee}$ for its smooth dual.

\begin{Prop}\label{Prop:Whit_lift}
    Let $\pi\in\mathrm{Irr}_{\epsilon}(\Gt)$ such that $\pi^{\s_{I_{+}}(I_{+})}\neq 0$. Then as vector spaces
    \begin{equation*}
        \mathrm{Hom}_{\C}(\pi_{(U^{-},\psi)},\C)\cong [(\pi^{\vee})^{K_{+}}]^{(U^{-}_{\kappa},\psi_{\kappa}^{-1})}.
    \end{equation*}
\end{Prop}

\begin{proof}
    By Frobenius reciprocity and passing to coinvariants we have
    \begin{equation*}
        \mathrm{Hom}_{\C}(\pi_{(U^{-},\psi)},\C)\cong \mathrm{Wh}_{\psi}(\pi).
    \end{equation*}
    By dualizing we get
    \begin{equation*}
        \mathrm{Wh}_{\psi}(\pi)\cong \mathrm{Hom}_{G}(\mathcal{V},\pi^{\vee}).
    \end{equation*}
    Since $\pi^{\s_{I_{+}}(I_{+})}\neq 0$ we may use that the pair $(\s_{I_{+}}(I_{+}),\mathrm{triv})$ is a Bushnell-Kutzko type \cite[Lemma 5.8]{GGK24} to get
    \begin{equation*}
        \mathrm{Hom}_{G}(\mathcal{V},\pi^{\vee})\cong \mathrm{Hom}_{\Hcal}(\mathcal{V}^{\s_{I_{+}}(I_{+})},(\pi^{\vee})^{\s_{I_{+}}(I_{+})}).
    \end{equation*}
    By \cite[Lemmas 5.1-5.3]{GGK24}, it follows that $\mathcal{W}^{\s_{I_{+}}(I_{+})}\cong \mathcal{V}^{\s_{I_{+}}(I_{+})}$ as $\Hcal$-modules. Thus 
    \begin{equation*}
        \mathrm{Hom}_{\Hcal}(\mathcal{V}^{\s_{I_{+}}(I_{+})},(\pi^{\vee})^{\s_{I_{+}}(I_{+})})\cong \mathrm{Hom}_{\Hcal}(\mathcal{W}^{\s_{I_{+}}(I_{+})},(\pi^{\vee})^{\s_{I_{+}}(I_{+})}).
    \end{equation*}
    Again since $(\s_{I_{+}}(I_{+}),\mathrm{triv})$ is a Bushnell-Kutzko type we have 
    \begin{equation*}
        \mathrm{Hom}_{\Hcal}(\mathcal{W}^{\s_{I_{+}}(I_{+})},(\pi^{\vee})^{\s_{I_{+}}(I_{+})})\cong \mathrm{Hom}_{G}(\mathcal{W},\pi^{\vee}).
    \end{equation*}
    By Frobenius reciprocity and because $\psi$ is a one dimensional representation we have
    \begin{equation*}
        \mathrm{Hom}_{G}(\mathcal{W},\pi^{\vee})\cong \mathrm{Hom}_{I^{-}_{+}}(\psi_{\kappa}^{-1},\pi^{\vee})\cong (\pi^{\vee})^{(I^{-}_{+},\psi_{\kappa}^{-1})}.
    \end{equation*}
    Finally we can compute the space of invariants in two steps because
    \begin{equation*}
        (\pi^{\vee})^{(I^{-}_{+},\psi_{\kappa}^{-1})}\cong [(\pi^{\vee})^{K_{+}}]^{(U^{-}_{\kappa},\psi_{\kappa}^{-1})}.
    \end{equation*}
\end{proof}

\section{Main Theorem}



\begin{Thm}\label{Thm:generalized_main}
    Let $\pi\in \mathrm{Irr}_{\epsilon}(\Gt)$. Let $\chi\in\mathrm{Hom}(T_{\kappa},\C^{\times})$, which can be inflated to a character of $\s_{I}(I)$. Let $\Ocal_{\chi}\subset \mathrm{Hom}(T_{\kappa},\C^{\times})$ be the $W\ltimes Y/\Yt$-orbit of $\chi$. Let $N$ be the number of free $W$-orbits in $\mathcal{O}_{\chi}$. Let $\psi\in \mathrm{Hom}(U^{-},\C^{\times})$ be a non-degenerate character of conductor $(\varpi)$.

    If $\pi^{(\s_{I}(I),\chi)}\neq 0$, then the Whittaker dimension with respect to $\psi$ is at least $N$. 
\end{Thm}

\begin{proof}
    Since $\pi^{(\s_{I}(I),\chi)}\neq 0$, Proposition \ref{Prop:Finite_Generic} implies that $\pi^{K_{+}}$ as a $G_{\kappa}$-representation contains $N$ non-isomorphic regular principal series. Thus it follows that for any non-degenerate character $\psi:U^{-}\rightarrow \C^{\times}$ we have 
    \begin{equation*}
        \mathrm{dim}(\pi^{K_{+}})^{(U^{-}_{\kappa},\psi_{\kappa}^{-1})}\geq N.
    \end{equation*}
    Thus by Proposition \ref{Prop:Whit_lift},
    \begin{equation*}
        \mathrm{dim}\pi_{(U^{-},\psi)}\geq N.
    \end{equation*}
\end{proof}

Now Theorem \ref{Thm:Main_Thm} follows from Theorem \ref{Thm:generalized_main} by setting $\chi=\chi_{\mathrm{triv}}$ to be the trivial character. In this case, the map $\phi: Y/\Yt\rightarrow \Ocal_{\chi_{\mathrm{triv}}}$ defined by $y+\Yt\mapsto \phi_{y}$ is a bijection, because $\s_{I} = \s_{K}|_{I}$ is a $W$-equivariant splitting of $I$. By Lemma \ref{Lem:Char_equi}, $\phi$ is $W$-equivariant. Thus the number of free $W$-orbits in $Y/\Yt$ is equal to the number of free $W$-orbits in $\Ocal_{\chi_{\mathrm{triv}}}$.

%Next we introduce a genuine character of $\It$. We begin with a general statement. If $H\subset G$ is subgroup and $\s:H\rightarrow \overline{H}$ is a splitting of \eqref{Eqn:CExt} over $H$, then the map $\mu_{n}\times H\cong \overline{H}$ defined by $(\zeta,h)\mapsto \zeta \s(h)$ is a group isomorphism. 

%The map $\mu_{n}\times H\rightarrow \C^{\times}$ defined by $(\zeta,h)\mapsto \epsilon(\zeta)$ is a group homomorphism and if we pullback to $\overline{H}$ via the isomorphism $\mu_{n}\times H\cong \overline{H}$ we get an $\epsilon$-genuine character. We will write $\phi_{\s}:\overline{H}\rightarrow \C^{\times}$ for this $\epsilon$-genuine character.

\newpage

\begin{center}{\bf Acknowledgments}\end{center}

This work began with a conversation with Dipendra Prasad at the IASM-BIRS 5-day workshop on the wave front set and endoscopic classification in Hangzhou, China in October 2026. Prasad prompted the author with the question: Is there a group theoretic explanation for the genericity of genuine Iwahori-spherical representation of high degree BD covers?

I would also like to thank Wee Teck Gan and Gordan Savin for comments on an earlier draft.

\bibliographystyle{amsalpha}
\bibliography{Hecke_algebra2_bio}


\end{document}

%%%%%%%%%%%%%%%%%%%%%%%%%%%%%%%%%%%%%%%%%%%%%%%
%%%%%%%%%%%%%%%%%%%%%%%%%%%%%%%%%%%%%%%%%%%%%%%
%%%%%%%%%%%%%%% TIKZCD Template code
%%%%%%%%%%%%%%%%%%%%%%%%%%%%%%%%%%%%%%%%%%%%%%%
%%%%%%%%%%%%%%%%%%%%%%%%%%%%%%%%%%%%%%%%%%%%%%%


\begin{equation*}
    \begin{tikzcd}
        A \ar[d, "{f}"]  \ar[r, "{g}"] & B \ar[d, "{h}"] \\
        C \ar[r, "{i}"] &  D
    \end{tikzcd}
\end{equation*}





\begin{comment}
Recall the natural isomorpihsm $T_{\kappa}\cong Y\otimes \kappa^{\times}$. Under this isomorphism and with our fixed choice of $\epsilon_{q-1}$, the character group of $T_{\kappa}$ can be identified with $X/(q-1)X$ in a $W$-equivariant way via $x(y\otimes u) = \epsilon(u^{\la x,y\ra})$, where $x\in X$, $y\in Y$ and $u\in\kappa^{\times}$.

Now for each $y\in Y$ we have the character $\phi_{y}|_{T_{0}}$. It is trivial on $T_{1}:=T\cap K_{+}$ and thus descends to a character of $T_{\kappa}$. We would like to identify this character as an element of $X/(q-1)X$.

The homomorphism $\mathscr{B}:Y\rightarrow X$ defined by $y\mapsto (y^{\prime}\mapsto B(y,y^{\prime}))$. This map is also $W$-equivariant.

Recall that for $y,z\in Y$ and $u\in \kappa^{\times}$ we have $\phi_{y}(z\otimes u) = (\varpi,u)_{n}^{B(y,z)}$. From this and the isomorphism $X/(q-1)X\cong X\otimes \widehat{\kappa^{\times}}$, we see that under the above identification $\widehat{T_{\kappa}}\cong X/(q-1)X$ the character $\phi_{y}$ corresponds to $\frac{q-1}{n}\mathscr{B}(y)+(q-1)X$.

Note that for any $m\in \Z$,
\begin{equation*}
   mX =\{x\in X|\la x,y\ra\in m\Z \text{, for all }y\in Y\}
\end{equation*}

Thus it follows that $\Yt =\mathscr{B}^{-1}((q-1)X)$.

This induces an injective $W$-equivariant group homomorphism 
\begin{equation*}
    \mathscr{B}_{n}:Y/\Yt\hookrightarrow (\frac{q-1}{n})X/(q-1)X.
\end{equation*}

The group $(\frac{q-1}{n})X/(q-1)X$ is a finite abelian group of size $n^{rk(\Gbb)}$, where $rk(\Gbb)$ is the rank of $\Gbb$.

Note that if $\Yt = nY$, then $\mathscr{B}_{n}:Y/\Yt\hookrightarrow (\frac{q-1}{n})X/(q-1)X$ is a $W$-equivariant isomorphism. 

The natural isomorphism $T_{0}/T_{1}\cong Y\otimes \kappa^{\times}$ combined with $\otimes-\mathrm{Hom}$ adjunction induces the following chain of isomorphisms.

\begin{align*}
    \mathrm{Hom}(T_{0}/T_{1},\C^{\times})\cong& \mathrm{Hom}(Y\otimes \kappa,\C^{\times})\\
    \cong& \mathrm{Hom}(Y,\hat{\kappa})\\
    \cong& X\otimes \hat{\kappa}\\
    \cong& X/(q-1)X.
\end{align*}

Under this chain of isomorphisms the element $\phi_{y}|_{T_{0}}$ (which is trivial on $T_{1}$) maps to $\mathscr{B}_{(q-1)}(y)+(q-1)X$.



Let $\mathscr{S}$ be the set of splittings of $I$. Note that for any $\s\in\mathscr{S}$, the genuine character $\phi_{\s}$ restricted to $\s_{I_{+}}(I_{+})$ is trivial. Thus the pair $(\It,\phi_{\s})$ is a depth $0$ Iwahori type.

The set $\mathscr{S}$ is a torsor under the group $\mathrm{Hom}(I,\mu_{n})$, where the action is given by point-wise multiplication.

Let $\mathrm{Hom}_{\epsilon}(\It,\mu_{n}(\C))\subset \mathrm{Hom}(\It,\mu_{n}(\C))$ be the subset of $\epsilon$-genuine characters. The set $\mathrm{Hom}_{\epsilon}(\It,\mu_{n}(\C))$ is a torsor under the group $\mathrm{Hom}(I,\mu_{n}(\C))$, where the action is given by point-wise multiplication (after pulling back along $\mathrm{pr}$). 

The map $\mathrm{Hom}(I,\mu_{n})\rightarrow \mathrm{Hom}(I,\mu_{n}(\C))$ defined by $f\mapsto \epsilon\circ f$ is an isomorphism of groups. 


\begin{Lem}
    The map $\mathscr{S}\rightarrow \mathrm{Hom}_{\epsilon}(\It,\mu_{n}(\C))$ defined by $s \mapsto \phi_{s}$ is an $\mathrm{Hom}(I,\mu_{n})\cong \mathrm{Hom}(I,\mu_{n}(\C))$-equivariant isomorphism.
\end{Lem}
\end{comment}